\honest{Arguing ``By Definition''}{Eliezer Yudkowsky}{2008}

I open with a plucked chicken which, having two legs and no feathers, is ``by definition'' a human. \nb{The joke is Diogenes' plucked chicken; see ``Similarity Clusters.''}

People who argue by definition usually start from visible features, find them in a dictionary, and conclude, ``Therefore, by definition, atheism is a religion!'' But features that both sides accept are rarely what is in dispute. If the listener doubted that Socrates has two legs, the reply would be ``Whaddaya mean Socrates is bipedal? That's what we're arguing about''; so the legs are probably not what is at stake.

You can legitimately go from Socrates's human shape to his vulnerability to hemlock. That inference is probable, not certain, and it rests on the world containing clusters of similar things, which no definition changes. A dictionary defining humans as ``all featherless bipeds except Socrates'' would not make him less like us. Argued correctly, it goes: Socrates has two arms and feet, speaks Greek, uses tools and has every property of Homo sapiens I can see, so I guess he has human biochemistry and, like every human tested, dies of hemlock.

Suppose I have seen Socrates with herbologists preparing an antidote. Then the argument has to move to the biochemistry of coniine. To insist that humans are mortal ``by definition'' is to throw away everything you know about Socrates except that he is human, like calling a coin 50\% heads after you have seen it land heads, or insisting that Frodo has ten fingers after seeing nine. Bayesian probability forbids refusing to condition on new evidence. Not every new fact matters: if Socrates has nine fingers, that does not noticeably change my estimate about hemlock, because losing a finger would not change the rest of his biochemistry, whatever the dictionary says about ten fingers. The legitimate inference rests on the clusters and on causal biology. Done right, it is said plainly: coniine in hemlock paralyzes the muscles and kills by asphyxiation; humans are vulnerable to hemlock.

``By definition'' comes out ``on exactly those occasions'' when the default inference has been called into doubt, meaning, in effect, forget what you heard about the herbologists. So too with ``X, by definition, is a Y!'', as in: atheists believe God does not exist, a negative belief is still a belief, so atheism answers theological questions, so atheism is by definition a religion. No one feels the need to say that Hinduism is a religion by definition; it is an actual religion. Atheism does not resemble the central members of the category, so people insist it is a religion by definition, ``because it isn't true any other way.'' They say it to sneak in a connotation. \nb{Where a definition is itself a rule that assigns consequences, as in law, the argument can be legitimate. In Kaufman v. McCaughtry (2005) an American appeals court wrote that ``Atheism is Kaufman's religion'' for the purposes of the First Amendment, and found that refusing an atheist prisoner's request to form a study group had violated his rights. No connotation was needed.}

I have kept track of this phrase for thirteen years, ``though not with literal statistics, I fear.'' By eyeballing, ``by definition'' outside mathematics is among the most alarming signals of flawed argument I have found, up there with ``Hitler,'' ``God,'' ``absolutely certain'' and ``can't prove that.'' \nb{That tally is the post's evidence for its claim about when people use the phrase. The cases the post analyzes, the herbologists and the atheism argument, are the author's own wording; no one is quoted.} The first correct use I spotted outside mathematics was by Feynman, and I have spotted more since; I do not show any of them. You are probably better off deleting it from your vocabulary, always when you are tempted to say it in italics or with an exclamation mark. ``That's a bad idea by definition!''
